\section{Variational foundations}\label{sec:foundations}

For completeness, this appendix supplies the two variational ingredients used above. The notation is that of Section~\ref{sec:cone}; all spaces may be complex unless real functions are specified. These arguments require neither a maximum principle nor uniqueness of an optimizer.


\begin{lemma}[Compactness of normalized near-optimizers]
\label{foundation:compact}
The sharp inequality \eqref{cone:gn} has an optimizer.
If $M(f_n)\to M_{c,a}$, $A_a(f_n)\to1$, and $\|f_n\|_p^p\to p/2$,
then after translations and passage to a subsequence,
$f_n$ converges strongly in $H^{m/2}$ to a normalized optimizer. After a constant phase, that optimizer is real. The family of normalized optimizers is compact modulo translations and constant phases.
\end{lemma}
\begin{proof}
The hypotheses make the sequence a normalized maximizing sequence. In particular, the same conclusion holds if $A_a(f_n)=1$ and $E_a(f_n)<0$, since the sharp inequality then forces $\|f_n\|_p^p\to p/2$.
Local Sobolev nonvanishing gives translations for which a subsequence
has a nonzero weak $H^{m/2}$ limit $Q$. Local Rellich compactness
also gives almost-everywhere convergence after a subsequence.
Write
\[
 m_1=M(Q)>0,\quad a_1=A_a(Q)>0,\quad
 m_2=M_{c,a}-m_1,\quad a_2=1-a_1.
\]
Hilbert norm splitting, the Brezis--Lieb lemma, and
\eqref{cone:gn} give
\[
 M_{c,a}^r\le m_1^ra_1+m_2^ra_2,\qquad a_1+a_2=1.
\]
If $0<m_1<M_{c,a}$, both coefficients $m_j^r$ are strictly below
$M_{c,a}^r$, which is impossible. Hence $m_1=M_{c,a}$, and the
same inequality forces $a_1=1$. Strong energy-space convergence
and equality in the sharp inequality follow.
To prove attainment, normalize an arbitrary maximizing sequence
by amplitude and dilation so that $M(f_n)=M_{c,a}$ and
$A_a(f_n)=1$. Its nonlinear norms converge to $p/2$, and the
identical splitting argument applies; negative energy is not
needed once that convergence is known. The same proof applies to a sequence of normalized optimizers. The constant-phase assertion follows from the next lemma.
\end{proof}

\begin{lemma}[Constant phase of optimizers]\label{foundation:constant-phase}
Every complex optimizer is a constant phase times a real function.
\end{lemma}
\begin{proof}
Write $f=A+iB$, where $A,B$ are real. Evenness of $a$ gives
$A_a(f)=A_a(A)+A_a(B)$. With
$C=(p/2)M_{c,a}^{-r}$, Minkowski, the sharp inequality on each real
component, and two-term H\"older yield
\begin{align*}
 \|f\|_p^2
 &\le\|A\|_p^2+\|B\|_p^2\\
 &\le C^{2/p}\left[
 A_a(A)^{2/p}M(A)^{1-2/p}
 +A_a(B)^{2/p}M(B)^{1-2/p}\right]\\
 &\le C^{2/p}A_a(f)^{2/p}M(f)^{1-2/p}.
\end{align*}
For an optimizer all inequalities are equalities. Strict convexity
of $L^{p/2}$ implies $B^2=cA^2$ almost everywhere, unless one
component vanishes. Repeat this equality argument after multiplying
$f$ by $e^{i\pi/4}$. If both signs of $B/A=\pm\sqrt c$ occurred
on sets of positive measure, the new squared components could
not be proportional. When $c=1$ the new components instead have
disjoint nonzero supports, again giving strict Minkowski inequality.
Thus only one sign is possible, and $f$ has constant phase.
\end{proof}

Let $Q$ be a real optimizer normalized by $M(Q)=M_{c,a}$ and
$A_a(Q)=1$. Its Euler--Lagrange equation is
\[
 a(D)Q+\omega_QQ-|Q|^{2r}Q=0,\qquad
 \omega_Q=r/M_{c,a}>0.
\]

\begin{lemma}[Strict phase coercivity]\label{foundation:phase}
The real quadratic form
\[
 L_-=a(D)+\omega_Q-|Q|^{2r}
\]
is nonnegative, has kernel $\operatorname{span}_{\mathbb R}\{Q\}$,
and obeys
\[
 \langle b,L_-b\rangle\ge c_Q\|b\|_{H^{m/2}}^2
 \quad\text{if }b\text{ is real and }\langle b,Q\rangle=0,
\]
for some $c_Q>0$.
\end{lemma}
\begin{proof}
Retain $A=A_a(Q)$ and $M=M(Q)$ in the algebra, so that
$\omega_Q=rA/M$. For real $b\in H^{m/2}$ set
$B=M(b)$ and $C_1=A_a(b)$.
The sharp inequality applied to $Q+i\sqrt t\,b$ gives
\begin{equation}\label{foundation:kernel-comparison}
 F(t):=\int(Q^2+tb^2)^{r+1}
 \le(r+1)(1+tB/M)^r(A+tC_1),\qquad t\ge0.
\end{equation}
The right derivatives at zero prove $L_-\ge0$.
If $b$ belongs to its form kernel, the first derivatives agree, and
\[
 \int |Q|^{2r}b^2=C_1+rAB/M.
\]
The pointwise second-order remainder for $F$ is nonnegative.
Comparison of second-order remainders in
\eqref{foundation:kernel-comparison} and Fatou's lemma yield
\[
 D:=\int |Q|^{2r-2}b^4
 \le(r-1)AB^2/M^2+2BC_1/M.
\]
On the zero set of $Q$, the integrand is interpreted through
the remainder: it is infinite if $r<1$ and $b\ne0$, equals $b^4$
if $r=1$, and equals zero if $r>1$.
Cauchy--Schwarz gives
\[
 D\ge\frac{(C_1+rAB/M)^2}{(r+1)A}.
\]
Combining the inequalities yields
$(C_1-AB/M)^2\le0$. Consequently all inequalities are equalities.
Equality in Cauchy--Schwarz implies
$b^2=(B/M)Q^2$ on $\{Q\ne0\}$. Its integral already equals $B$,
so $b=0$ on the remaining zero set as well.
The form kernel is a linear space containing $Q$. Apply the same
conclusion to $b+Q$ and subtract the identity for $b^2$.
It follows that $b$ is one constant multiple of $Q$.

Finally, multiplication by $|Q|^{2r}$ is compact from $H^{m/2}$
to $H^{-m/2}$. Indeed $|Q|^{2r}\in L^{p/(2r)}$,
H\"older with two $L^p$ factors makes the tail small, and the
embedding into local $L^p$ is compact because $p$ lies strictly
below the critical Sobolev exponent. Thus $L_-$ is a compact
form perturbation of the coercive form $a(D)+\omega_Q$.
Nonnegativity and its one-dimensional kernel give the asserted
coercivity on the $L^2$ orthogonal complement of $Q$.
\end{proof}

